\documentclass[11pt]{amsart}

\usepackage[T1]{fontenc}
\usepackage[utf8]{inputenc}
\usepackage{amsmath,amssymb,amsthm}
\usepackage[margin=1in]{geometry}
\usepackage{array}
\usepackage{hyperref}

\newcommand{\OO}{\mathbb{O}}
\newcommand{\RR}{\mathbb{R}}
\newcommand{\QQ}{\mathbb{Q}}
\newcommand{\ZZ}{\mathbb{Z}}
\newcommand{\FF}{\mathbb{F}}
\newcommand{\Nm}{\mathrm{N}}
\newcommand{\cj}[1]{\overline{#1}}

\numberwithin{equation}{section}

\theoremstyle{definition}
\newtheorem{definition}{Definition}[section]
\theoremstyle{plain}
\newtheorem{theorem}[definition]{Theorem}
\newtheorem{proposition}[definition]{Proposition}
\theoremstyle{remark}
\newtheorem*{tracing}{Sources and tracing}

\title[How far Johannes Kirmse got in 1924]{How far Johannes Kirmse got in
1924:\\ a verified exposition of his arithmetic\\ of the integral octonions}
\author{Claude Opus 4.7}
\address{Anthropic}
\date{Preliminary version -- 20 May 2026}

\begin{document}
\maketitle

\begin{abstract}
This is a self-contained exposition of Johannes Kirmse's 1924 paper on the
arithmetic of the octonions~\cite{Kirmse1924}, written by the AI agent at the
direction of Jens K\"oplinger as part of the research record of the
repository~\cite{repo}; the verification it reports was carried out under
prompt~100 of that record.  Its single purpose is to establish, from
Kirmse's original text and by independent computation, \emph{exactly} what
Kirmse achieved and where his one substantive error lies.

The conclusion is this.  Kirmse built the arithmetic theory of the
integral octonions in full: the units and the multiplication table,
conjugation and norm, the composition identities, modules, orders and
maximal orders, an ideal theory, and from it the classical eight-squares
theorem.  His one substantive error is a miscount: he states
(\cite{Kirmse1924}, p.~70) that there are \emph{eight} maximal orders
containing the integer lattice~$J_0$, whereas there are \emph{seven}; and the
one maximal order he writes down explicitly, his~$J_1$, is itself \emph{not}
closed under multiplication.  Every formula below is traced to Kirmse's
original; each section closes with a \emph{Sources and tracing} remark.

This is a preliminary version.  The companion historical questions -- who
corrected the count, in what form, and where the error first entered the
later literature -- require Coxeter's 1946 paper and are deliberately left
untouched here.  An appendix records, as explicit speculation and at the
suggestion of J.~K\"oplinger, one possible human provenance of the error.
\end{abstract}

% ===============================================================
\section{Purpose and method}
% ===============================================================

Johannes Kirmse's name survives in two phrases.  To workers in
non-associative algebra it survives in \emph{Kirmse's identities}; to those
concerned with integral quadratic forms it survives in \emph{Kirmse's
mistake}.  The second phrase has all but erased the first, and with it the
fact that Kirmse's short 1924 paper is one of the founding documents of the
arithmetic of the octonions.

The present note was written to answer one question precisely: how far did
Kirmse get, and where exactly is the mistake?  It follows the methodology of
the repository~\cite{repo} in which it was produced: no claim -- Kirmse's,
or anyone's -- is taken on trust; every assertion is either traced to
Kirmse's printed text or verified from scratch by exact computation.  The
two are kept strictly separate.  Sections~\ref{sec:arith}--\ref{sec:after}
expound Kirmse's paper and trace every formula to it.  Section~\ref{sec:verify}
reports an independent computation, owing nothing to Kirmse or to any later
author, that tests his central claim.  Section~\ref{sec:assess} draws the
balance.

We write $\OO$ for the real algebra Kirmse studies -- the Graves--Cayley
octonions -- and reserve the symbols $J_0$, $J_1$ for the two lattices he
names.

% ===============================================================
\section{The elementary arithmetic of the Cayley numbers}\label{sec:arith}
% ===============================================================

Kirmse begins with an algebra of \emph{eight} units
$i_0, i_1, \ldots, i_7$, with $i_0 = 1$ the identity, and a real algebra
\[
  a = \alpha_0 + \alpha_1 i_1 + \alpha_2 i_2 + \cdots + \alpha_7 i_7,
  \qquad \alpha_0,\ldots,\alpha_7 \in \RR .
\]
The multiplication of the units is fixed by his table~(1).  Reading that
table, the seven products that carry a positive sign organise into seven
oriented triples
\begin{equation}\label{eq:triples}
  (1,2,3),\ (1,5,4),\ (1,6,7),\ (2,4,7),\ (2,6,5),\ (3,4,6),\ (3,5,7),
\end{equation}
with the rule: for each triple $(a,b,c)$,
\[
  i_a i_b = i_c,\quad i_b i_c = i_a,\quad i_c i_a = i_b
  \qquad\text{(and the reverse products carry a minus sign),}
\]
together with $i_j^2 = -1$ for $j=1,\ldots,7$ and $i_0 i_j = i_j i_0 = i_j$.
Written out, the products of the seven imaginary units are

\smallskip
\[
\renewcommand{\arraystretch}{1.25}
\begin{array}{c|ccccccc}
\cdot & i_1 & i_2 & i_3 & i_4 & i_5 & i_6 & i_7\\\hline
i_1 & -1   & i_3  & -i_2 & -i_5 & i_4  & i_7  & -i_6\\
i_2 & -i_3 & -1   & i_1  & i_7  & -i_6 & i_5  & -i_4\\
i_3 & i_2  & -i_1 & -1   & i_6  & i_7  & -i_4 & -i_5\\
i_4 & i_5  & -i_7 & -i_6 & -1   & -i_1 & i_3  & i_2\\
i_5 & -i_4 & i_6  & -i_7 & i_1  & -1   & -i_2 & i_3\\
i_6 & -i_7 & -i_5 & i_4  & -i_3 & i_2  & -1   & i_1\\
i_7 & i_6  & i_4  & i_5  & -i_2 & -i_3 & -i_1 & -1
\end{array}
\]
\smallskip

\noindent
This is Kirmse's table~(1) in full.  He notes that the problem he set
himself -- an eight-unit system containing the quaternions and carrying the
fourth power of the sum of squares as a norm -- admits ``several, but not
essentially different'' solutions, of which~(1) is one.

To each $a$ Kirmse associates the \emph{symmetric} number (the conjugate)
and the \emph{Relativnorm}
\begin{align}
  \cj a &= \alpha_0 - \alpha_1 i_1 - \cdots - \alpha_7 i_7,
  \label{eq:conj}\\
  \Nm(a) &= a\cj a = \sum_{j=0}^{7} \alpha_j^{\,2}.
  \label{eq:norm}
\end{align}
Conjugation reverses products,
\begin{equation}\label{eq:antiauto}
  \cj{a b} = \cj b\,\cj a .
\end{equation}
Developing the eight products $a, a i_1, \ldots, a i_7$ in the units yields
an $8\times8$ matrix $M[a]$ of a linear substitution; its determinant, the
\emph{Grundnorm}, equals $\Nm(a)^4$, and from the structure of $M[ab]$
Kirmse obtains the composition law
\begin{equation}\label{eq:comp}
  \Nm(ab) = \Nm(a)\,\Nm(b).
\end{equation}
Calling $a$ \emph{regular} when $\Nm(a)\neq 0$, he solves the one-sided
equations $ax=c$ and $ya=c$:
\begin{equation}\label{eq:div}
  x = \frac{\cj a\,c}{\Nm(a)},
  \qquad
  y = \frac{c\,\cj a}{\Nm(a)} .
\end{equation}
Specialising $c=ab$ gives the two \emph{composition identities}
\begin{equation}\label{eq:kirmse-id}
  \cj a\,(a b) = \Nm(a)\,b,
  \qquad
  (a b)\,\cj b = \Nm(b)\,a,
\end{equation}
so that ``for products of three factors in which two consecutive factors are
conjugate, the associative law holds'' -- in particular
$(\cj a a)b = \cj a(ab)$ and $a(b\cj b) = (ab)\cj b$.  Finally Kirmse calls
$a_1,a_2,a_3$ \emph{coplanar} when corresponding imaginary components are
proportional; coplanar numbers commute and associate, every power $a^m$ is
coplanar with $a$, and ``coplanar numbers behave like algebraic numbers.''

\begin{tracing}
The eight units and the multiplication table reproduce Kirmse's table~(1),
\cite{Kirmse1924} p.~64; we have rendered his printed array of $64$ products
$i_j i_k$ as the seven oriented triples~\eqref{eq:triples} plus
$i_j^2=-1$ and $i_0 i_k=i_k$, and the $7\times7$ array above is that same
table written out (it was checked entry-by-entry against the scan of p.~64).
The conjugate~\eqref{eq:conj} and norm~\eqref{eq:norm} are Kirmse's, p.~64
(he names $\cj a$ the number ``symmetric to~$a$'' and $\Nm(a)$ the
\emph{Relativnorm}).  Equation~\eqref{eq:antiauto} is Kirmse's~(2), p.~64.
The matrix $M[a]$ is his~(3) and the product expansion his~(4), p.~65; the
composition law~\eqref{eq:comp} is his~(5), p.~66, derived there from the
column structure of $M[ab]$.  The division formulas~\eqref{eq:div} are
his~(7\,a) and~(7\,b), pp.~66--67; the composition
identities~\eqref{eq:kirmse-id} are the content of his~(8), p.~67.  The
remarks on coplanar numbers follow his~(9)--(10), p.~67.  All of this is
Kirmse's \S1, ``Die elementare Arithmetik der Cayleyschen Zahlen.''
\end{tracing}

% ===============================================================
\section{Modules, orders, and the count of maximal orders}\label{sec:orders}
% ===============================================================

Kirmse now restricts to the rational octonions $K$ -- those $a$ with all
$\alpha_j\in\QQ$ -- and develops a theory of lattices in $K$.

\begin{definition}
A \emph{module} is a subset of $K$ closed under integer linear combinations.
It is \emph{finite} if it consists of all $\gamma_0 a_0 + \cdots + \gamma_7 a_7$
($\gamma_j\in\ZZ$) for eight independent numbers $a_0,\ldots,a_7$.  A finite
module is an \emph{Integrit\"atsbereich} -- an \emph{order} -- if with any
two numbers $a,b$ it also contains the products $ab$ and $ba$.  An order is a
\emph{most-comprehensive order} (\emph{umfassendster Integrit\"atsbereich})
-- a \emph{maximal order} -- if it contains
\begin{equation}\label{eq:J0}
  J_0 = [\,1, i_1, i_2, i_3, i_4, i_5, i_6, i_7\,]
\end{equation}
and cannot be enlarged by adjoining a further number of $K$ without ceasing
to be an order.
\end{definition}

The lattice $J_0$ of octonions with integer components is itself an order.
Kirmse's structural theorem about \emph{any} order is obtained as follows.
If $a$ lies in an order with $\ZZ$-basis $a_0,\ldots,a_7$, the products
$a a_j$ have integer coordinates in that basis; eliminating the basis from
the eight resulting linear relations gives an integer determinant $D[a]$
with
\begin{equation}\label{eq:degree8}
  D[a] = 0,
\end{equation}
so every number of an order satisfies a monic integral equation of degree~8.
Because conjugates and coplanar combinations again satisfy such equations,
$a+\cj a$ and $a\cj a$ are rational integers; and multiplying by a suitable
unit moves any chosen component to the front.  Kirmse concludes:

\begin{theorem}[Kirmse, p.~69]\label{thm:doubles}
The double of every number belonging to an order $J\subseteq K$ has
exclusively integer components.
\end{theorem}

Thus the numbers of an order have components that are integers or
half-integers, and the maximal orders may be sought by a finite search among
lattices between $J_0$ and $\tfrac12 J_0$.  At this point Kirmse writes
(\cite{Kirmse1924}, p.~70), verbatim:

\begin{quote}\itshape
``Dieser Satz erlaubt sofort, durch eine einfache endliche Rechnung die
s\"amtlichen umfassendsten Integrit\"atsbereiche \"uber $J_0$ aufzusuchen.
\textnormal{\textbf{Man findet acht solche}}, deren jeder die Gestalt hat''
\end{quote}
\begin{equation}\label{eq:eight}
  J = \Bigl[\, J_0,\ \tfrac12\bigl(i_{\nu_1}+i_{\nu_2}+i_{\nu_3}+i_{\nu_4}\bigr),\
  \tfrac12\textstyle\sum_{k=0}^{7} i_k \,\Bigr],
\end{equation}
the indices $\nu_1,\nu_2,\nu_3,\nu_4$ all distinct, drawn from
$\{0,1,\ldots,7\}$.  That is: \emph{Kirmse asserts there are eight maximal
orders.}  He adds that each contains exactly fourteen elements of the
``middle kind'' $\tfrac12(\text{four units})$, falling into seven
complementary pairs.

He then exhibits \emph{one} of the eight explicitly.  Calling it $J_1$,
\begin{equation}\label{eq:J1}
  J_1 = \bigl[\, J_0,\ a_1,\ a_2,\ a_3,\ a_4 \,\bigr],
\end{equation}
\begin{equation}\label{eq:J1gens}
\begin{aligned}
  a_1 &= \tfrac12\bigl(1 + i_1 + i_2 + i_3\bigr), &
  a_2 &= \tfrac12\bigl(i_4 + i_5 + i_6 + i_7\bigr),\\
  a_3 &= \tfrac12\bigl(1 + i_1 + i_4 + i_5\bigr), &
  a_4 &= \tfrac12\bigl(1 + i_2 + i_4 + i_7\bigr),
\end{aligned}
\end{equation}
with $J_0$ of index $\Nm(J_1)=16$ in $J_1$.  Every later section of the
paper works in this $J_1$.

\begin{tracing}
The definitions of module, finite module, order (\emph{Integrit\"atsbereich})
and maximal order (\emph{umfassendster Integrit\"atsbereich}), and the
lattice $J_0$ of~\eqref{eq:J0}, are Kirmse's, \cite{Kirmse1924} p.~68.  The
integral degree-8 equation~\eqref{eq:degree8} is his~(11), p.~69, derived in
his \S2 by the elimination just summarised.  Theorem~\ref{thm:doubles} is the
italicised \emph{Satz} at the foot of his p.~69.  The claim of eight maximal
orders and the schematic form~\eqref{eq:eight} are on his p.~70; the quoted
German is verbatim, and the emphasis on ``acht'' (\emph{eight}) is added.
The order $J_1$ and its generators~\eqref{eq:J1}--\eqref{eq:J1gens}, and the
index $16$, are Kirmse's, p.~70.  All of this is his \S2, ``Moduln und
Integrit\"atsbereiche im Gebiete der Cayleyschen Zahlen mit rationalen
Komponenten.''
\end{tracing}

% ===============================================================
\section{What Kirmse built on $J_1$}\label{sec:after}
% ===============================================================

The remaining three sections, all set in $J_1$, show the reach of the
construction; we summarise them only to record how far Kirmse carried the
theory.

In \S3 (``Die ungeraden Ideale\ldots'') he develops one- and two-sided ideals,
their norms, and primitive and primary ideals.  In \S4 (``Die Einheiten und
die geraden Ideale in $J_1$'') he finds the units: \emph{sixteen} in $J_0$
(the $\pm i_j$) and \emph{two hundred and forty} in $J_1$ --
the $\pm i_j$ together with $224$ half-integer units -- and observes that
they ``form a group in an extended sense, in which all group-theoretic
theorems hold except the associative law and its consequences,'' i.e.\ a
Moufang loop; he counts $2160 = 240\cdot 9$ numbers of norm~2.  In \S5 (``Die
Primzerlegung\ldots'') he gives the prime decomposition, finds that every
ideal of $J_1$ (and every odd ideal of $J_0$) is principal, computes the
ideal zeta function
\begin{equation}\label{eq:zeta}
  Z(s) = \zeta(4s)\,\zeta(4s-3),
\end{equation}
and derives the eight-squares theorem: the number of representations of an
odd $n>0$ as a sum of eight integer squares is
\begin{equation}\label{eq:eightsquares}
  r_8(n) = 16\sum_{d\mid n} d^{\,3}.
\end{equation}
A postscript adds a Minkowski geometry-of-numbers proof that all ideals are
principal.

These results -- the $240$ units, \eqref{eq:zeta}, \eqref{eq:eightsquares}
-- are correct.  They are facts about the $E_8$ lattice as a quadratic form
and, as Section~\ref{sec:verify} will make clear, do not depend on the one
claim that fails.

\begin{tracing}
The section titles and contents are Kirmse's \S\S3--5, \cite{Kirmse1924}
pp.~70--82: the ideal theory of \S3 (from p.~70); the units and the
count~$240$, the ``group in an extended sense,'' and $2160=240\cdot9$ in \S4,
p.~76; the zeta function~\eqref{eq:zeta} (his $Z(s)$, p.~77) and the
eight-squares theorem~\eqref{eq:eightsquares} (his numbered conclusions~2--3,
p.~78); the Minkowski argument is the \emph{Nachschrift}, p.~82.  Equations
\eqref{eq:zeta}--\eqref{eq:eightsquares} restate his results in his own
notation.
\end{tracing}

% ===============================================================
\section{Independent verification}\label{sec:verify}
% ===============================================================

Kirmse asserts~\eqref{eq:eight} -- eight maximal orders -- ``durch eine
einfache endliche Rechnung,'' but prints no part of that computation.  We
carry it out, in exact arithmetic, from his data alone: the multiplication
table~(1) of Section~\ref{sec:arith} and the order $J_1$
of~\eqref{eq:J1}--\eqref{eq:J1gens}.  Nothing from any later author enters.

\medskip\noindent\textbf{(a) The table is sound.}  Built from the seven
triples~\eqref{eq:triples}, the multiplication satisfies the composition
law~\eqref{eq:comp}, $\Nm(uv)=\Nm(u)\Nm(v)$, on every tested pair.  Kirmse's
$\OO$ is a genuine (composition) octonion algebra.

\medskip\noindent\textbf{(b) $J_1$ is not closed.}  $J_1$ is the
$\ZZ$-span of the twelve numbers $1,i_1,\ldots,i_7,a_1,a_2,a_3,a_4$.  By
bilinearity, $J_1$ is an order if and only if all $144$ products of these
generators lie in $J_1$.  Ten of them do not.  The simplest witness uses two
of Kirmse's own generators:
\begin{equation}\label{eq:witness}
  a_1\cdot a_3
  = \tfrac12(1{+}i_1{+}i_2{+}i_3)\cdot\tfrac12(1{+}i_1{+}i_4{+}i_5)
  = \tfrac12\bigl(i_1+i_2+i_4+i_7\bigr).
\end{equation}
The number $\tfrac12(i_1+i_2+i_4+i_7)$ does not belong to $J_1$: its
half-integer pattern is not one of the sixteen that $J_1$ admits.  Hence
\[
  J_1\cdot J_1 \not\subseteq J_1 ;
\]
\emph{$J_1$ is not closed under multiplication, and so is not an order at
all.}

\medskip\noindent\textbf{(c) There are seven maximal orders, not eight.}  A
maximal order containing $J_0$ is a unimodular even lattice, hence
corresponds to a doubly-even self-dual binary code of length~$8$; there are
exactly $30$ such codes (the coordinate relabellings of the extended Hamming
code).  Testing each of the $30$ candidate lattices for closure under
table~(1) leaves exactly
\[
  \boxed{\ \text{seven}\ }
\]
that are closed.  Kirmse's $J_1$ is among the $30$ candidates but is one of
the $23$ that fail.  His count~\eqref{eq:eight} is therefore too large by
one, and the representative he chose was one of the non-orders.

\medskip\noindent\textbf{(d) The structure of the error.}  Closure of a
candidate lattice is the conjunction of two independent sub-conditions: that
every product of a unit with a half-integer element again lies in the
lattice, and that every product of two half-integer elements does.  The
first sub-condition alone -- invariance under multiplication by the units
-- is satisfied by exactly \emph{eight} of the thirty candidates; imposing
the second removes exactly one.  The eight unit-invariant lattices are
precisely the seven genuine maximal orders together with Kirmse's~$J_1$, and
$J_1$ is the unique one of the eight that is not an order.  Concordantly,
every one of $J_1$'s closure failures is a product of two half-integer
elements -- ten of the sixteen products $a_i\cdot a_j$ -- while all $128$
products involving a unit close.  The count ``eight'' is thus not arbitrary:
it is exactly the number delivered by the first, weaker sub-condition.

\begin{tracing}
This section traces not to Kirmse but to an independent computation.  The
multiplication is that of Section~\ref{sec:arith} (Kirmse's table~(1));
the lattice $J_1$ is~\eqref{eq:J1}--\eqref{eq:J1gens} (Kirmse's p.~70).  The
arithmetic is exact and is reproducible from the scripts
\texttt{verify\_kirmse\_1924.py} and \texttt{verify\_kirmse\_1924\_forensic.py}
(in \texttt{python\_project/src/}) of the repository~\cite{repo}; the runs
were performed under prompts~100 and~103 of the research record.  Item~(b)'s
witness~\eqref{eq:witness} was also checked by hand.  The count of $30$
doubly-even self-dual binary codes of length~$8$ is classical; items~(c)
and~(d) re-derive it, and the figures ``seven'' and ``eight,'' the identity
of the eight, and the block in which $J_1$ fails are all outputs of exact
computation, not quotations.
\end{tracing}

% ===============================================================
\section{Assessment}\label{sec:assess}
% ===============================================================

What Kirmse accomplished in these twenty pages is considerable, and is set
out in Sections~\ref{sec:arith}--\ref{sec:after}: a working octonion algebra;
conjugation, norm, and the composition law; the composition
identities~\eqref{eq:kirmse-id} that bear his name; the notions of order and
maximal order; the structural Theorem~\ref{thm:doubles}; an ideal theory; the
$240$-unit Moufang loop; and, as the paper's announced goal, a derivation of
the eight-squares theorem~\eqref{eq:eightsquares}.  Much of the modern
arithmetic of the integral octonions is already here, in 1924.

His one substantive error is exactly located.  It is not in the algebra, nor
in the ideal theory, nor in the final theorem -- all of which concern the
$E_8$ lattice as a quadratic form and stand.  It is the single claim,
\eqref{eq:eight}, that there are \emph{eight} maximal orders over $J_0$.
There are seven.  And the one maximal order Kirmse wrote down and used,
$J_1$, is itself not multiplicatively closed: of the orders he actually
exhibited, the score is none out of one.  This is the precise content of
what later became known as ``Kirmse's mistake.''  It does not diminish the
rest; it is one miscount in a paper that otherwise built a theory.

One matter is deliberately left open here: who detected the error, and in
what form the count was corrected.  That is the subject of Coxeter's 1946
paper, and this exposition will be revised once that paper has been read from
the source in the same manner as Kirmse's.

A second matter is now largely accounted for.  Section~\ref{sec:verify}(d)
shows that ``eight'' is exactly the number of candidate lattices invariant
under multiplication by the units, and that those eight are the seven genuine
maximal orders together with~$J_1$.  Since Kirmse printed no computation, the
identification of his ``simple finite computation'' with this unit-invariance
test is a reconstruction, not a quotation -- but it reproduces both his
count and his choice of~$J_1$ exactly.  What is \emph{not} recoverable from
his text is the route by which Kirmse himself arrived there.  A frankly
speculative reconstruction of that route, contributed by J.~K\"oplinger, is
set out in Appendix~\ref{app:provenance}; it is kept separate precisely
because it is speculation, not record.

A small biographical note, since it differs from some secondary accounts:
the paper's by-line reads ``Johannes Kirmse \emph{in Hannover},'' it was
presented to the Saxon Academy by G.~Herglotz in the session of 5~May 1924,
and it is signed ``Hannover, am 15.~April 1924.''

\begin{tracing}
The list of accomplishments refers back to
Sections~\ref{sec:arith}--\ref{sec:after} and through them to Kirmse's
\S\S1--5.  The localisation of the error refers to~\eqref{eq:eight}
(Kirmse, p.~70) and to Section~\ref{sec:verify}.  The by-line, the
presentation by Herglotz, and the place and date of signature are read from
\cite{Kirmse1924}, pp.~63 and~82.  The deferred questions concern
H.~S.~M.~Coxeter, \emph{Integral Cayley numbers}, Duke Math.\ J.~\textbf{13}
(1946), 561--578, not consulted for this preliminary version.
\end{tracing}

% ===============================================================
\appendix
\section{A speculative note on provenance}\label{app:provenance}
% ===============================================================

\emph{This appendix differs in kind from the body.  Everything above is
either traced to Kirmse's printed text or verified by computation; what
follows is neither.  It is an openly speculative reconstruction -- proposed
by Jens K\"oplinger -- of how a human being might have been led to the error
of Section~\ref{sec:orders}.  It is placed in an appendix, and marked
plainly, so that it cannot be mistaken for the record.}

\medskip

Kirmse's paper states its general claim first and only afterwards descends to
the particular order~$J_1$.  Human inquiry often runs the other way: a
concrete instance, found and checked, comes first, and the general statement
is distilled from it.  One may therefore picture -- and it is only a
picture -- Kirmse proceeding in that order: a worked example in hand; a
generalisation drawn from it to the form ``$J_0$, half-integer middle
elements, $\tfrac12\sum i_k$''; a count of the resulting class; and a
representative chosen from it.  On this reading the misstep enters at the
generalisation.  Not in the \emph{form}, which is correct -- every genuine
maximal order has it -- but in the leap from a single instance to a complete
classification, asserted on a criterion left unspoken.
Section~\ref{sec:verify}(d) gives that unspoken criterion a definite shape:
closure under the units, a condition necessary for an order but not
sufficient, which admits eight lattices where seven are orders.  The
representative Kirmse then exhibited, $J_1$, would on this reading have been
drawn from a class delimited by a test he never stated -- and which does
not, in fact, delimit the orders.

What the paper permits, and what it does not, must be said plainly.  It does
\emph{not} permit the story above to be told as history.  Kirmse exhibits
exactly one concrete realisation, $J_1$, and $J_1$ is not an order; if a
working first example ever existed, it has left no trace.  The order of
\emph{exposition} -- general, then particular -- is not evidence of the
order of \emph{discovery}.  What is certain is only the logical content of
Section~\ref{sec:verify}: closure is a two-part test, and the second part was
not carried through.  The reconstruction above is consistent with that
content, and is offered for no stronger reason than consistency.

Why record a speculation at all.  The anatomy of an error -- how it is
made, how it survives a careful reader, how it is at last detected and
corrected -- is itself worth study, and a candid hypothesis about the human
step, marked unmistakably as hypothesis, can serve that study better than
silence.  That is the spirit in which this appendix is offered, and the
reason it is quarantined from the body.

\begin{tracing}
This appendix traces to no source.  The speculative reconstruction of a human
provenance was contributed by Jens K\"oplinger; the present author has
rendered it in a form consistent with the verified findings of
Sections~\ref{sec:verify}--\ref{sec:assess} and with the rule that
speculation be labelled as such.  Its only factual anchors -- the two-part
closure test, the eight unit-invariant lattices, and the block in which
$J_1$ fails -- are those of Section~\ref{sec:verify}(d).
\end{tracing}

% ===============================================================
\begin{thebibliography}{9}

\bibitem{Kirmse1924}
J.~Kirmse,
\emph{\"Uber die Darstellbarkeit nat\"urlicher ganzer Zahlen als Summen von
acht Quadraten und \"uber ein mit diesem Problem zusammenh\"angendes
nichtkommutatives und nichtassoziatives Zahlensystem},
Berichte \"uber die Verhandlungen der S\"achsischen Akademie der
Wissenschaften zu Leipzig, Mathematisch-Physische Klasse \textbf{76} (1924),
63--82.

\bibitem{repo}
J.~K\"oplinger,
\emph{An order on the Leech lattice from a $\ZZ_3$-symmetric triple-octonion
product: research repository},
\url{https://bitbucket.org/jenskoeplinger/leech_alg}, 2026.
Verification scripts: \texttt{python\_project/src/verify\_kirmse\_1924.py}
and \texttt{verify\_kirmse\_1924\_forensic.py};
research record: prompts~100, 103, 108.

\end{thebibliography}

\end{document}
